\lesson{6}{Jan 05 2022 Wed (20:02:45)}{Reactions in Our World}{Unit 4}

In \bf{Chemical Reactions}, the loss or gain of \bf{Electrons} does not change
an element's identity. But during a \bf{Nuclear Reaction}, a \bf{Nucleus} is
split apart or $2$ smaller \bf{Nuclei} are \un{combined}.

This changes the identity of the element because \bf{protons} and \bf{neutrons}
are gained or lost by the atom over the course of the \bf{Nuclear Reaction}.
These \it{"Lost"} \bf{protons} and \bf{Neutrons} don't just disappear, they
become new \un{particles}.


\subsubsection{Fission and Fusion}

Albert Einstein is one of the most famous scientists in history. Most people are
familiar with one of his most famous accomplishments, the formula:

\begin{align}
  E = mc^2
.\end{align}

This formula shows a relationship between \it{mass} and \it{energy}.

In regular \bf{Chemical Reactions}, \it{mass} and \it{energy} are both
conserved.

This means the measured amounts of \it{mass} and \it{energy} are the same at the
beginning and end of a chemical reaction. In a nuclear reaction, scientists have
found that \it{mass} can be converted to \it{energy}. This is important,
especially when you see that a small loss of \it{mass} from an atom's nucleus
releases extremely large amounts of energy!

Elements that have a greater mass than \bf{Iron} are \un{Unstable}, and are
often \bf{\un{\it{Radioactive}}}. These large atoms may undergo one of the
several types of nuclear reactions. In a process called \it{Nuclear Fission}, an
atom of \bf{Uranium} may break apart into two smaller atoms. When this happens,
the resulting atoms are actually less massive in total than the original
\bf{Uranium} atoms.

The "missing" mass was converted to energy according to Einstein's formula E =
mc2. Multiplying the lost mass by the speed of light squared determines the
amount of energy given off by each uranium atom as it splits.

It is possible to make the nuclei of two hydrogen atoms join together to form
one helium nucleus. This requires that the nuclei be hurled at each other at a
very high speed. This process occurs in the sun, but it can possibly be
replicated on Earth with the use of lasers or magnets, or in the center of an
atomic bomb. The process is called nuclear fusion.

When the two nuclei combine, the new nucleus has less mass than the two single,
separate nuclei. When the nuclei are forced together, this extra mass is
released as energy. The amount of energy released can be predicted using
Einstein's famous formula E = mc2.

\newpage

